\ifdefined\gkver\else\def\gkver{student}\fi
\documentclass[\gkver]{gaokaozhenti}
\begin{document}

\chapter{2003年全国理综新课程}
%% paperType: 理综物理部分

\begin{enumerate}
\item
%% number: 16
%% paperName: 2003年全国理综新课程第 16 题 6 分
%% typeId: 1
%% type: 单选题
%% chapter: 电场
%% point: 库仑定律
%% method: 无
%% score: 6
%% degree: 650
%% duplicateId: 0
%% body:
如图所示，三个完全相同的金属小球 a、b、c 位于等边三角形的三个顶点上。a 和 c 带正电，b 带负电，a 所带电量的大小比 b 的小。已知 c 受到 a 和 b 的静电力的合力可用图中四条有向线段中的一条来表示，它应是\xzanswer{B}

\onepicture[w=4.8cm]{\includesvg{figs/16}}

\fourchoices[answer=B]
{$F_{1}$}
{$F_{2}$}
{$F_{3}$}
{$F_{4}$}

%% answer: B
%% memo:
\memoanswer{
分析 c 球，c 球受 a 球的斥力和 b 球的引力的方向可知，它们的合力在 ac 延长线与 cb 区域范围内，故 C、D 错误；因 $q_{a}<q_{b}$，则 $F_{bc}>F_{ac}$，所以表示 $F_{bc}$ 的矢量线比表示 $F_{ac}$ 的矢量线长，其合力靠近 $F_{bc}$，故 $F_{2}$ 才是符合实际的，故 A 错误，B 正确。
}

\item
%% number: 17
%% paperName: 2003年全国理综新课程第 17 题 6 分
%% typeId: 1
%% type: 单选题
%% chapter: 原子和原子核
%% point: 人工核反应
%% method: 无
%% score: 6
%% degree: 750
%% duplicateId: 0
%% body:
下面列出的是一些核反应方程

\ce{^{30}_{15}P -> ^{30}_{14}Si + X}，\ce{^{9}_{4}Be + ^{2}_{1}H -> ^{10}_{5}B + Y}，\ce{^{4}_{2}He + ^{4}_{2}He -> ^{7}_{3}Li + Z}

其中\xzanswer{D}

\fourchoices[answer=D]
{X 是质子，Y 是中子，Z 是正电子}
{X 是正电子，Y 是质子，Z 是中子}
{X 是中子，Y 是正电子，Z 是质子}
{X 是正电子，Y 是中子，Z 是质子}

%% answer: D
%% memo:
\memoanswer{
核反应方程满足质量数守恒、核电荷数守恒。在 \ce{^{30}_{15}P -> ^{30}_{14}Si + X} 方程中质量数没变，核电荷数少 1，所以 X 为正电子；在 \ce{^{9}_{4}Be + ^{2}_{1}H -> ^{10}_{5}B + Y} 方程中，核电荷数没变，质量数少 1，所以 Y 为中子；在 \ce{^{4}_{2}He + ^{4}_{2}He -> ^{7}_{3}Li + Z} 方程中，核电荷数少 1，质量数少 1，所以 Z 为质子，故 D 正确。
}

\item
%% number: 18
%% paperName: 2003年全国理综新课程第 18 题 6 分
%% typeId: 1
%% type: 单选题
%% chapter: 机械波
%% point: 波长频率波速
%% method: 无
%% score: 6
%% degree: 650
%% duplicateId: 0
%% body:
简谐机械波在给定的媒质中传播时，下列说法中正确的是\xzanswer{D}

\fourchoices[answer=D]
{振幅越大，则波传播的速度越快}
{振幅越大，则波传播的速度越慢}
{在一个周期内，振动质元走过的路程等于一个波长}
{振动的频率越高，则波传播一个波长的距离所用的时间越短}

%% answer: D
%% memo:
\memoanswer{
简谐机械波在给定的媒质中传播时，速度由介质决定，与振幅无关，故 A、B 错误；一个周期内，振动质元运动的路程为四个振幅，与波长无关，故 C 错误；波传播一个波长的距离所用的时间等于一个周期，振动的频率越高，周期越短，波传播一个波长的距离所用的时间越短，故 D 正确。
}

\item
%% number: 19
%% paperName: 2003年全国理综新课程第 19 题 6 分
%% typeId: 1
%% type: 单选题
%% chapter: 力物体的平衡
%% point: 物体系平衡
%% method: 无
%% score: 6
%% degree: 550
%% duplicateId: 0
%% body:
如图所示，一个半球形的碗放在桌面上，碗口水平，$O$ 点为其球心，碗的内表面及碗口是光滑的。一根细线跨在碗口上，线的两端分别系有质量为 $m_{1}$ 和 $m_{2}$ 的小球，当它们处于平衡状态时，质量为 $m_{1}$ 的小球与 $O$ 点的连线与水平线的夹角为 $\alpha=60^{\circ}$。两小球的质量比 $\frac{m_{2}}{m_{1}}$ 为\xzanswer{A}

\onepicture[w=5.0cm]{figs/19.png}

\fourchoices[answer=A]
{$\frac{\sqrt{3}}{3}$}
{$\frac{\sqrt{2}}{3}$}
{$\frac{\sqrt{3}}{2}$}
{$\frac{\sqrt{2}}{2}$}

%% answer: A
%% memo:
\memoanswer{
若采用整体法分析，因为绳与碗口处的弹力未知无法得解，故应该用隔离法分析。对 $m_{2}$ 进行受力分析可知 $T=m_{2}g$；对 $m_{1}$ 进行受力分析，如图所示，由几何关系得 $\theta=\frac{\alpha}{2}$，由共点力平衡条件得 $T\cos\frac{\alpha}{2}+N\cos\frac{\alpha}{2}=m_{1}g$，$T\cos\alpha=N\cos\alpha$，联立解得 $\frac{m_{2}}{m_{1}}=\frac{\sqrt{3}}{3}$，故 A 正确。

\onepicture[w=2.2cm]{figs/19_an.jpg}
}

\item
%% number: 20
%% paperName: 2003年全国理综新课程第 20 题 6 分
%% typeId: 1
%% type: 单选题
%% chapter: 气体性质
%% point: 热力学第一定律
%% method: 无
%% score: 6
%% degree: 650
%% duplicateId: 0
%% body:
如图所示，固定容器及可动活塞 P 都是绝热的，中间有一导热的固定隔板 B，B 的两边分别盛有气体甲和乙。现将活塞 P 缓慢地向 B 移动一段距离，已知气体的温度随其内能的增加而升高，则在移动 P 的过程中，\xzanswer{C}

\onepicture[w=5.0cm]{figs/20.png}

\fourchoices[answer=C]
{外力对乙做功；甲的内能不变}
{外力对乙做功；乙的内能不变}
{乙传递热量给甲；乙的内能增加}
{乙的内能增加；甲的内能不变}

%% answer: C
%% memo:
\memoanswer{
当活塞 P 向 B 移动时，外力做功，气体乙的内能增加，气体乙的温度升高，乙通过导热板 B 向气体甲传热，因此甲的内能增加，故 A、B、D 错误，C 正确。
}

\item
%% number: 21
%% paperName: 2003年全国理综新课程第 21 题 6 分
%% typeId: 1
%% type: 单选题
%% chapter: 光学
%% point: 光电效应
%% method: 无
%% score: 6
%% degree: 550
%% duplicateId: 0
%% body:
如图，当电键 K 断开时，用光子能量为 $2.5\UeV$ 的一束光照射阴极 P，发现电流表读数不为零。合上电键，调节滑线变阻器，发现当电压表读数小于 $0.60\UV$ 时，电流表读数仍不为零；当电压表读数大于或等于 $0.60\UV$ 时，电流表读数为零。由此可知阴极材料的逸出功为\xzanswer{A}

\onepicture[w=4.8cm]{figs/21.png}

\fourchoices[answer=A]
{$1.9\UeV$}
{$0.6\UeV$}
{$2.5\UeV$}
{$3.1\UeV$}

%% answer: A
%% memo:
\memoanswer{
由光电效应方程 $E_{k}=h\nu-W_{0}$，又 $h\nu=2.5\UeV$，得 $W_{0}=h\nu-E_{k}=2.5\UeV-0.60\UeV=1.9\UeV$，故 A 正确。
}

\item
%% number: 22
%% paperName: 2003年全国理综新课程第 22 题 6 分
%% typeId: 1
%% type: 单选题
%% chapter: 磁场
%% point: 带电粒子在磁场中的运动
%% method: 无
%% score: 6
%% degree: 450
%% duplicateId: 0
%% body:
$K^{-}$ 介子衰变的方程为：$K^{-}\to\pi^{-}+\pi^{0}$，其中 $K^{-}$ 介子和 $\pi^{-}$ 介子带负的基元电荷，$\pi^{0}$ 介子不带电。一个 $K^{-}$ 介子沿垂直于磁场的方向射入匀强磁场中，其轨迹为圆弧 AP，衰变后产生的 $\pi^{-}$ 介子的轨迹为圆弧 PB，两轨迹在 P 点相切，它们的半径 $R_{K^{-}}$ 与 $R_{\pi^{-}}$ 之比为 $2:1$。$\pi^{0}$ 介子的轨迹未画出。由此可知 $\pi^{-}$ 的动量大小与 $\pi^{0}$ 的动量大小之比为\xzanswer{C}

\onepicture[w=4.8cm]{figs/22.png}

\fourchoices[answer=C]
{$1:1$}
{$1:2$}
{$1:3$}
{$1:6$}

%% answer: C
%% memo:
\memoanswer{
带电粒子在匀强磁场中做匀速圆周运动的半径公式 $r=\frac{mv}{qB}$，$K^{-}$ 介子和 $\pi^{-}$ 介子电荷量又相同，说明它们的动量大小之比是 $2:1$，方向相反；由动量守恒得 $\pi^{0}$ 介子的动量大小是 $\pi^{-}$ 介子的三倍，方向与 $\pi^{-}$ 介子的速度方向相反，故 C 正确。
}

\item
%% number: 23
%% paperName: 2003年全国理综新课程第 23 题 15 分
%% typeId: 4
%% type: 实验题
%% chapter: 电路
%% point: 伏安法测电阻
%% method: 无
%% score: 15
%% degree: 550
%% duplicateId: 0
%% body:
用伏安法测量电阻阻值 $R$，并求出电阻率 $\rho$。

给定电压表（内阻约为 $50\UkO$）、电流表（内阻约为 $40\UO$）、滑线变阻器、电源、电键、待测电阻（约为 $250\UO$）及导线若干。
\begin{enumerate}
	\item
	画出测量 $R$ 的电路图。
	\item
	图 1 中的 6 个点表示实验中测得的 6 组电流 $I$、电压 $U$ 的值，试写出根据此图求 $R$ 值的步骤：\tkanswer[5cm]{①作 $U-I$ 直线，舍去左起第 2 点，其余 5 个点尽量靠近直线且均匀分布在直线两侧；②求该直线的斜率 $k$，则 $R=k$}。求出的电阻值 $R=$ \tkanswer[2.2cm]{$229\UO$}。（保留 3 位有效数字）
	\item
	待测电阻是一均匀材料制成的圆柱体，用游标为 50 分度的卡尺测量其长度与直径，结果分别如图 2、图 3 所示。由图可知其长度为 \tkanswer[2.0cm]{$0.800\Ucm$}，直径为 \tkanswer[2.0cm]{$0.194\Ucm$}。
	\item
	由以上数据可求出 $\rho=$ \tkanswer[2.5cm]{$8.46\times10^{-2}\UOm$}。（保留 3 位有效数字）
\end{enumerate}

\twopicture[showlabel=false, wA=4.6cm, wB=4.4cm, align=right]{figs/23a.png}{figs/23b.png}

%% answer:
\jdanswer{
\begin{enumerate}
	\item
	电路图如图 1 或图 2 所示。
	\item
	①作 $U-I$ 直线，舍去左起第 2 点，其余 5 个点尽量靠近直线且均匀分布在直线两侧；②求该直线的斜率 $k$，则 $R=k$。$R=229\UO$（$221\sim237\UO$ 均为正确）。
	\item
	$0.800\Ucm$，$0.194\Ucm$。
	\item
	$8.46\times10^{-2}\UOm$（$8.16\times10^{-2}\sim8.76\times10^{-2}\UOm$ 均为正确）。
\end{enumerate}
}
%% memo:
\memoanswer{
由题意设计电路图，因没有明确的说明，所以分压、限流电路都可以。对于电流表的接法，由于待测电阻（约 $250\UO$）与电流表内阻（约 $40\UO$）相近，远小于电压表内阻（$50\UkO$），且满足 $R_{x}<\sqrt{R_{A}R_{V}}$，所以应采用外接法，电路图如图 1 或图 2 所示。

\twopicture[showlabel=false, wA=4.0cm, wB=3.8cm]{figs/23_an1.png}{figs/23_an2.png}

（2）作出 $U-I$ 图线，舍去左起第 2 点，其余 5 个点在一条直线上，不在这条线上的点尽量均匀分布在直线两侧；求该直线的斜率 $k$，则 $R=k=229\UO$（$221\sim237\UO$ 均为正确）。

（3）因为游标为 50 分度，所以游标卡尺的精确度为 $\frac{1}{50}\Umm=0.02\Umm$，另外游标卡尺不需要估读，读出待测电阻的长度为 $0.800\Ucm$，直径为 $1\Umm+47\times0.02\Umm=1.94\Umm=0.194\Ucm$。

（4）将数据代入 $\rho=\frac{RS}{l}=R\frac{\pi d^{2}}{4l}$，得 $\rho=8.46\times10^{-2}\UOm$。
}

\item
%% number: 24
%% paperName: 2003年全国理综新课程第 24 题 15 分
%% typeId: 6
%% type: 计算题
%% chapter: 曲线运动
%% point: 人造地球卫星
%% method: 无
%% score: 15
%% degree: 650
%% duplicateId: 0
%% body:
中子星是恒星演化过程的一种可能结果，它的密度很大。现有一中子星，观测到它的自转周期为 $T=\frac{1}{30}\Us$。问该中子星的最小密度应是多少才能维持该星体的稳定，不致因自转而瓦解。计算时星体可视为均匀球体。（引力常数 $G=6.67\times10^{-11}\Umq/(\Ukg\cdot\Usq)$）

%% answer:
\jdanswer{
$\rho=1.27\times10^{14}\Ukgmc$
}
%% memo:
\memoanswer{
考虑中子星赤道处一小块物质，只有当它受到的万有引力大于或等于它随星体一起旋转所需的向心力时，中子星才不会瓦解。

设中子星的密度为 $\rho$，质量为 $M$，半径为 $R$，自转角速度为 $\omega$，位于赤道处的小块物质质量为 $m$，则有

$\frac{GMm}{R^{2}}=m\omega^{2}R$，$\omega=\frac{2\pi}{T}$，$M=\rho\cdot\frac{4}{3}\pi R^{3}$

由以上各式得：$\rho=\frac{3\pi}{GT^{2}}$

代入数据解得：$\rho=1.27\times10^{14}\Ukgmc$。
}

\item
%% number: 25
%% paperName: 2003年全国理综新课程第 25 题 18 分
%% typeId: 6
%% type: 计算题
%% chapter: 电磁感应
%% point: 双杆问题
%% method: 无
%% score: 18
%% degree: 450
%% duplicateId: 0
%% body:
两根平行的金属导轨，固定在同一水平面上，磁感强度 $B=0.50\UT$ 的匀强磁场与导轨所在平面垂直，导轨的电阻很小，可忽略不计。导轨间的距离 $l=0.20\Um$。两根质量均为 $m=0.10\Ukg$ 的平行金属杆甲、乙可在导轨上无摩擦地滑动，滑动过程中与导轨保持垂直，每根金属杆的电阻为 $R=0.50\UO$。在 $t=0$ 时刻，两杆都处于静止状态。现有一与导轨平行，大小为 $0.20\UN$ 的恒力 $F$ 作用于金属杆甲上，使金属杆在导轨上滑动。经过 $t=5.0\Us$，金属杆甲的加速度为 $a=1.37\Umsq$，问此时两金属杆的速度各为多少？

\onepicture[w=6.5cm, align=right]{figs/25.png}

%% answer:
\jdanswer{
$v_{1}=8.15\Ums$，$v_{2}=1.85\Ums$
}
%% memo:
\memoanswer{
设任一时刻 $t$ 两金属杆甲、乙之间的距离为 $x$，速度分别为 $v_{1}$ 和 $v_{2}$，经过很短时间 $\Delta t$，杆甲移动距离 $v_{1}\Delta t$，杆乙移动距离 $v_{2}\Delta t$，回路面积改变

$\Delta S=\left[(x-v_{2}\Delta t)+v_{1}\Delta t\right]l-lx=(v_{1}-v_{2})l\Delta t$

由法拉第电磁感应定律，回路中的感应电动势

$E=B\frac{\Delta S}{\Delta t}$

回路中的电流

$I=\frac{E}{2R}$

杆甲的运动方程

$F-BlI=ma$

由于作用于杆甲和杆乙的安培力总是大小相等、方向相反，所以两杆的动量（$t=0$ 时为 0）等于外力 $F$ 的冲量

$Ft=mv_{1}+mv_{2}$

联立以上各式解得

$v_{1}=\frac{1}{2}\left[\frac{Ft}{m}+\frac{2R}{B^{2}l^{2}}(F-ma)\right]$

$v_{2}=\frac{1}{2}\left[\frac{Ft}{m}-\frac{2R}{B^{2}l^{2}}(F-ma)\right]$

代入数据得

$v_{1}=8.15\Ums$，$v_{2}=1.85\Ums$。
}

\item
%% number: 34
%% paperName: 2003年全国理综新课程第 34 题 22 分
%% typeId: 6
%% type: 计算题
%% chapter: 运动和力
%% point: 传送带
%% method: 无
%% score: 22
%% degree: 250
%% duplicateId: 0
%% body:
一传送带装置示意如图，其中传送带经过 AB 区域时是水平的，经过 BC 区域时变为圆弧形（圆弧由光滑模板形成，未画出），经过 CD 区域时是倾斜的，AB 和 CD 都与 BC 相切。现将大量的质量均为 $m$ 的小货箱一个一个在 A 处放到传送带上，放置时初速为零，经传送带运送到 D 处，D 和 A 的高度差为 $h$。稳定工作时传送带速度不变，CD 段上各箱等距排列，相邻两箱的距离为 $L$。每个箱子在 A 处投放后，在到达 B 之前已经相对于传送带静止，且以后也不再滑动（忽略经 BC 段时的微小滑动）。已知在一段相当长的时间 $T$ 内，共运送小货箱的数目为 $N$。这装置由电动机带动，传送带与轮子间无相对滑动，不计轮轴处的摩擦。求电动机的平均输出功率 $\overline{P}$。

\onepicture[w=8.0cm, align=right]{figs/34.png}

%% answer:
\jdanswer{
$\overline{P}=\frac{Nm}{T}\left(\frac{N^{2}L^{2}}{T^{2}}+gh\right)$
}
%% memo:
\memoanswer{
以地面为参考系（下同），设传送带的运动速度为 $v_{0}$，在水平段运输的过程中，小货箱先在滑动摩擦力作用下做匀加速运动，设这段路程为 $s$，所用时间为 $t$，加速度为 $a$，则对小箱有

$s=\frac{1}{2}at^{2}$ ①

$v_{0}=at$ ②

在这段时间内，传送带运动的路程为

$s_{0}=v_{0}t$ ③

由以上可得

$s_{0}=2s$ ④

用 $f$ 表示小箱与传送带之间的滑动摩擦力，则传送带对小箱做功为

$W=fs=\frac{1}{2}mv_{0}^{2}$ ⑤

传送带克服小箱对它的摩擦力做功

$W_{0}=fs_{0}=2\cdot\frac{1}{2}mv_{0}^{2}$ ⑥

两者之差就是克服摩擦力做功发出的热量

$Q=\frac{1}{2}mv_{0}^{2}$ ⑦

可见，在小箱加速运动过程中，小箱获得的动能与发热量相等。

$T$ 时间内，电动机输出的功为

$W=\overline{P}T$ ⑧

此功用于增加小箱的动能、势能以及克服摩擦力发热，即

$W=\frac{1}{2}Nmv_{0}^{2}+Nmgh+NQ$ ⑨

已知相邻两小箱的距离为 $L$，所以

$v_{0}T=NL$ ⑩

联立⑦⑧⑨⑩，得

$\overline{P}=\frac{Nm}{T}\left(\frac{N^{2}L^{2}}{T^{2}}+gh\right)$
}

\end{enumerate}

\end{document}
